The door to the melted hall had been broken for as long as anybody’s grandmother could remember, which was rather the whole reason it was worth breaking into in the first place, because a door that hung crooked on one hinge and no longer properly latched was not, strictly, forbidden, whatever your mother said about the place.

“You said it wasn’t locked.”

“It’s not locked. It’s stuck. That’s different.” Tam put her shoulder against the old wood and shoved, and it gave with a groan that echoed considerably further than she’d expected, which she pretended not to notice, because Robbin was watching to see if she’d flinch and she was not, under any circumstances that morning, going to give him the satisfaction. “There. See? Not locked.”

“It’s dark in there.”

“It’s daytime, Pip. It’s dark like a room’s dark, not dark like anything’s going to eat you.”

“Something ate a whole dragon in there once,” Pip said, with the flat certainty of a boy of seven repeating something he’d heard and had no particular reason to doubt. “Granny Nell told me.”

“Granny Nell tells you a great many things,” Robbin said, “and about half of them are true, and I’ve never once worked out which half before it’s too late to matter.”

\sceneBreak

They went in single file, because the doorway wasn’t wide enough for three abreast, Tam first because it had been her idea and she’d sooner die than let either of the boys claim the going-in-first as their own, Robbin second because he’d rather be second than admit he’d wanted to be third, and Pip last, one hand fisted hard in the back of Robbin’s shirt, which nobody remarked on, because everybody present understood that a boy of seven was allowed exactly this much cowardice before it started counting against him properly.

The hall opened up bigger than any of them had pictured it, even having heard it described a hundred times over by a hundred different mouths across their whole short lives, and for a moment none of the three of them said anything at all, because there is a particular kind of quiet that a very large empty room makes a person fall into, regardless of how loud they’d been outside it a moment before.

Light came down through a gap in the roof on the eastern side, where nobody had ever quite finished the mending of it, and dust hung in the shaft of it, turning slow, and at the far end of the hall, half in shadow and half caught in that same grey light, sat the thing they had actually come to see.

\sceneBreak

“That’s it,” Tam breathed. “That’s the actual chair.”

“Doesn’t look like a chair. Looks like somebody melted a chair.”

“That’s because somebody \textit{did} melt it, stupid. That’s the whole point of it.” Tam started forward, and the boys followed, because whatever fear the doorway had cost them had apparently used itself up entirely on the getting-in, and none of them had any left over for the walking-across. “My granddad told my mam and my mam told me. There was a war. The last war, the real one, before all the little ones after it that don’t properly count. And there were dragons — three of them, or maybe more, nobody’s certain any longer — and one of the dragons breathed fire right over this whole chair because the queen sitting on it had gone wicked and mad, wickeder than anybody in a hundred years, and burning the chair was the only way anybody could think of to stop her being queen of it.”

“That’s not what my granny says.” Robbin folded his arms, which was his way of announcing that a proper argument was about to begin and everybody ought to settle in for it. “My granny says it wasn’t the queen at all who went mad. She says it was a \textit{king}, and that he was already dead, only nobody had told him yet, and he came walking up out of the ground with a whole army of dead men behind him, and the dragons burned him up right here in this room, and that’s why the floor’s still black in that one corner over there, see it, that’s not shadow, that’s \textit{burn}.”

“There’s no army of dead men in any of it. That’s a story for scaring babies with. I’m not a baby.”

“You believed the dragon story same as I believed the dead king one, so I don’t see you’ve got much room to talk.”

\sceneBreak

Pip, who had wandered a little apart from the other two while they argued, because he had worked out early in his short life that the surest way to avoid being dragged into somebody else’s fight was to look very interested in something else entirely, made a small, delighted sound from over by the base of the wreckage.

“I found a coin.”

“Let’s see it.”

He held it up, small and green-black with age, worn nearly smooth, and Tam took it and turned it over twice in the grey light, frowning at the shape still faintly visible on its face — a spiral, cut shallow, worn almost to nothing by whatever hands had passed it along before it ended up lost in the rubble of a room nobody official had swept in longer than any of their parents had been alive.

“That’s not proper coin. No king’s face on it. No house mark either.”

“Maybe it’s foreign,” Robbin offered.

“Maybe it’s just old,” Tam said, and pocketed it, because a coin with no face on it was still a coin, and a coin was always worth having about you, whatever else it was or wasn’t. None of them thought to wonder, and none of them would ever have any particular reason to, what exactly it had once been minted to prove, or whose hands had first decided that a government with nothing behind it but a promise needed a shape simple enough that anybody, anywhere, might trust it without being told why.

\sceneBreak

“There’s a story about a woman, too,” Robbin said, once the coin business had settled, in the lower voice a child uses when he’s about to say something he half wants to frighten his friends with and half wants to frighten himself. “My granny told me this one proper, not the baby version. There was a woman, back then, who could take anybody’s face. Not a mask. Her \textit{own} face, gone, and somebody else’s put on in its place, so good you’d never know unless you’d known the real one your whole life and looked very close. Granny says she still walks about even now, somewhere, in some city or other, wearing whoever she needs to be that week, and that if a stranger’s ever too kind to you too fast, that’s how you’d know.”

“That’s not real.”

“I never said it was real. I said it’s the story.” Robbin shrugged, with the particular satisfaction of a boy who has successfully unsettled his audience without having to commit to believing the thing himself. “Granny says nobody ever found out what happened to her proper. She just went. One day she was somewhere and the next she wasn’t anywhere anybody could point to, and that’s rather worse than a dead-and-gone story, because at least a dead-and-gone story finishes.”

Tam thought about that for a moment, standing in the grey light with a stranger’s worn-smooth coin in her pocket, and found that she agreed, though she could not have said properly why, that a story with no ending to it was a stranger and somehow larger thing to carry than a story that simply stopped.

\sceneBreak

“What about him, then,” Pip asked, pointing, and the other two turned to look at what he’d found — a mark on the wall near the base of the great melted ruin, old, half worn away by years of other hands touching the same spot for reasons none of them had ever bothered to ask their elders about, a rough shape that might once have been a hand, and beneath it, in letters gone shallow and soft with age, something that might once have said \textit{HE STOOD HERE}, though the last two letters had worn down to almost nothing at all.

“That’s him,” Tam said, with total, unearned confidence. “The man who killed the dead king with his own bare hand. Granny Nell told me that one too — he touched the dead thing and it just fell apart, turned to dust right there, because his hand was blessed, on account of he’d been born of fire himself, half a dragon, though he never told anybody that part until the very end.”

“That’s not what happened at all,” Robbin said, though with rather less conviction than his last correction, because he found, standing in front of the mark himself, that he did not, in fact, have his own competing version ready, and had to settle for simply disagreeing in principle rather than in any particular detail. “I don’t know what happened. But it wasn’t that.”

“How do you know it wasn’t, if you don’t know what it was instead?”

Robbin did not have an answer for that, and admitted as much with a shrug, and Pip, who had lost interest in the argument the moment it stopped being about anything he could hold in his hand, wandered off toward the fallen beam at the eastern side, where the last of the old timber had long since been carried away for firewood by people with more pressing uses for good wood than a museum nobody official had ever decided to make of this room, leaving only scorch marks and a wide gap of grey sky where a roof had once been.

\sceneBreak

“I’m hungry,” Pip announced, from across the hall, in the manner of a boy who had decided the whole expedition had run its course now that the coin and the wall-mark had both been properly investigated and there was nothing left in the room that could be eaten, worn, or kept.

“We only just got here.”

“I don’t care. I’m hungry, and Mam said I wasn’t to be late for the midday bread, and if I’m late again she’ll not let me come next time, and then you two will have to find the coin next time all on your own, and I don’t believe either of you is half so good at finding things as I am.”

Tam considered arguing this point, weighed it against her own stomach, which had been making its own opinions known for some while now that she’d stopped to notice, and decided the argument wasn’t worth the losing of a meal over.

“Fine. But I’m keeping the coin.”

“You can keep the coin. I don’t want it. It’s not even proper money.” Pip was already heading for the door, considerably less cautious about the going-out than he had been about the coming-in, the way boys of seven generally were about anything standing between themselves and bread. “Besides, Granny Nell says there’s a whole crate of old letters and things up in her attic, from when she was a girl my age, and she’s promised to show me one of these days, and I’d rather see those than an old chair, if I’m honest, because a chair’s just a chair, however melted, and a letter’s got somebody’s actual hand in it.”

\sceneBreak

They went out the way they’d come, single file, into a morning considerably brighter than the grey light they’d left behind, and the city around them was loud and ordinary in the way a city is loud and ordinary on any morning when nothing in particular is happening in it — a woman arguing with a fishmonger over the morning’s price, a cart wheel needing grease and complaining about it the whole length of the street, two old men arguing, with real feeling, about whether the coming winter would be a hard one, in the manner of old men who have argued about the weather every year of their lives and fully intend to go on doing so.

Nobody in that street looked twice at three children coming out of a doorway that had been broken for longer than any of them could properly remember, because there was nothing remarkable in it, and there had not been for a great many years now, and there was nothing in the whole of that ordinary morning to suggest that the room behind that door held anything a person ought to trouble themselves over, beyond a broken old chair, and a coin worth nothing in particular, and a mark on a wall that meant whatever the last person telling the story that morning had decided it ought to mean.

The throne sat where it had always sat, half-melted, entirely silent, indifferent to whichever version of itself the city told, and would go on sitting there long after Tam had lost the coin somewhere and forgotten she’d ever had it, and long after Robbin’s granny’s story about the faceless woman had been handed down once more, worn a little smoother in the handling, the way every good story wears smooth, the true edges of it rounding off a little further with each mouth that carried it forward.

And the city went on around it, the way cities do, entirely unbothered, because a broken chair in an empty room is only ever exactly as important as the people still living choose to make it, and on that particular bright, ordinary morning, with bread waiting and a whole day of considerably more interesting things ahead of them, the people still living had, without any of them quite deciding to, chosen not very important at all.